\subsection{Stokes 问题的正则性}
\label{sec:chap4-stokes-regularity-proof}

\renewcommand{\thestatement}{6.\arabic{statement}}
\renewcommand{\theHstatement}{ch04.sec06.\arabic{statement}}
\setcounter{statement}{0}
\newcounter{chapterfoursectionsixremark}
\renewcommand{\thechapterfoursectionsixremark}{6.\arabic{chapterfoursectionsixremark}}

\subsubsection{第一阶正则性}
\label{subsec:chap4-stokes-first-regularity}

本节的目标是证明 $k=0$ 时的 Theorem~\ref{thm:chap4-stokes-regularity}. 我们在 Section~\ref{subsec:chap4-stokes-higher-regularity} 中使用归纳论证证明 Theorem~\ref{thm:chap4-stokes-regularity} 的一般形式.

\paragraph{局部正则性}
\label{subsubsec:chap4-stokes-local-regularity}

证明 Theorem~\ref{thm:chap4-stokes-regularity} 的第一步是建立一个局部正则性结果. 注意, 这个结果不要求区域 $\Omega$ 满足任何假设. 此外, 它与解所满足的边界条件无关.

\begin{Theorem}
\label{thm:chap4-stokes-local-regularity}
设 $\Omega$ 是 $\mathbb R^d$ 的任意开子集, $f\in(L_{\mathrm{loc}}^2(\Omega))^d$, 且 $g\in H_{\mathrm{loc}}^1(\Omega)$. 对任意满足
\[
 (v,p)\in(H_{\mathrm{loc}}^1(\Omega))^d\times L_{\mathrm{loc}}^2(\Omega)
\]
并作为下列问题之解的 $(v,p)$:
\[
 \left\{
 \begin{aligned}
  -\Delta v+\nabla p&=f,&&\text{于 }\Omega,\\
  \operatorname{div}v&=g,&&\text{于 }\Omega,
 \end{aligned}
 \right.
\]
都有
\[
 (v,p)\in(H_{\mathrm{loc}}^2(\Omega))^d\times H_{\mathrm{loc}}^1(\Omega).
\]
更准确地说, 对任意满足 $\omega\subset\omega_2\subset\overline{\omega_2}\subset\Omega$ 的有界开集 $\omega,\omega_2$, 存在 $C(\omega,\omega_2)>0$, 使得
\[
 \|v\|_{H^2(\omega)}+\|p\|_{H^1(\omega)}
 \leq C(\omega,\omega_2)\bigl(\|f\|_{L^2(\omega_2)}+\|g\|_{H^1(\omega_2)}
 +\|v\|_{H^1(\omega_2)}+\|p\|_{L^2(\omega_2)}\bigr).
\]
\end{Theorem}

\begin{Proof}
由 Proposition~\ref{prop:chap3-local-sobolev-characterization}, 只需证明对任意 $\varphi\in\mathcal D(\Omega)$, 有 $\overline{\varphi v}\in(H^2(\mathbb R^d))^d$ 且 $\overline{\varphi p}\in H^1(\mathbb R^d)$, 并且 $\|\overline{\varphi v}\|_{H^2(\mathbb R^d)}+\|\overline{\varphi p}\|_{H^1(\mathbb R^d)}$ 具有适当估计.

为此, 令 $w=\overline{\varphi v}$, $q=\overline{\varphi p}$. 由于 $\varphi$ 光滑, 按照 Theorem~\ref{thm:chap3-laplace-dirichlet-regularity} 的证明中的同样计算, 立即得到 $w$ 和 $q$ 在 $\mathbb R^d$ 的分布意义下满足
\[
 \left\{
 \begin{aligned}
  -\Delta w+\nabla q&=\overline{\varphi f}-\overline{(\Delta\varphi)v}
  -2\overline{\nabla v\cdot\nabla\varphi}+\overline{p\nabla\varphi},\\
  \operatorname{div}w&=\overline{\varphi g}+\overline{v\cdot\nabla\varphi}.
 \end{aligned}
 \right.
\]
若将这些方程的第二端分别记为 $\widetilde f$ 和 $\widetilde g$, 则显然 $\widetilde f\in(L^2(\mathbb R^d))^d$, $\widetilde g\in H^1(\mathbb R^d)$, 且
\begin{align}
 \|\widetilde f\|_{L^2(\mathbb R^d)}
 &\leq \|\varphi\|_{L^\infty}\|f\|_{L^2(\omega_2)}
 +\|\Delta\varphi\|_{L^\infty}\|v\|_{L^2(\omega_2)}
 +2\|\nabla\varphi\|_{L^\infty}\|v\|_{H^1(\omega_2)}
 +\|\nabla\varphi\|_{L^\infty}\|p\|_{L^2(\omega_2)}
 \notag\\
 &\leq C_1(\varphi)\bigl(\|f\|_{L^2(\omega_2)}+\|g\|_{H^1(\omega_2)}
 +\|v\|_{H^1(\omega_2)}+\|p\|_{L^2(\omega_2)}\bigr),
 \label{eq:chap4-localized-stokes-f-bound}
\end{align}
以及
\begin{equation}
 \|\widetilde g\|_{H^1(\mathbb R^d)}
 \leq C_2(\varphi)\bigl(\|f\|_{L^2(\omega_2)}+\|g\|_{H^1(\omega_2)}
 +\|v\|_{H^1(\omega_2)}+\|p\|_{L^2(\omega_2)}\bigr).
 \label{eq:chap4-localized-stokes-g-bound}
\end{equation}
此外, 已有 Theorem~\ref{thm:chap4-stokes-homogeneous-wellposedness} 给出的估计, 即
\begin{equation}
 \|w\|_{H^1(\mathbb R^d)}
 \leq(\|\varphi\|_{L^\infty}+\|\nabla\varphi\|_{L^\infty})\|v\|_{H^1(\omega_2)}
 \leq C(\varphi)\|v\|_{H^1(\omega_2)},
 \label{eq:chap4-localized-stokes-w-bound}
\end{equation}
以及
\begin{equation}
 \|q\|_{L^2(\mathbb R^d)}\leq\|\varphi\|_{L^\infty}\|p\|_{L^2}.
 \label{eq:chap4-localized-stokes-q-bound}
\end{equation}
因此, 我们需要研究全空间 $\mathbb R^d$ 中的 Stokes 方程组
\begin{equation}
 \left\{
 \begin{aligned}
  -\Delta w+\nabla q&=\widetilde f,\\
  \operatorname{div}w&=\widetilde g.
 \end{aligned}
 \right.
 \label{eq:chap4-whole-space-localized-stokes}
\end{equation}
注意, 该方程组中出现的所有函数在紧集 $\overline\Omega$ 外均为零, 这使下文的所有运算都有根据.

我们用平移算子刻画 Sobolev 空间 $H^{k+1}(\mathbb R^d)$, 这种刻画由 Proposition~\ref{prop:chap3-sobolev-difference-quotient-characterization} 给出. 使用 Chapter~\ref{chap:sobolev-spaces} 的 Section~\ref{sec:chap3-sobolev-lipschitz-domains} 中引入的记号. 需要估计 $\delta_h w$ 和 $\delta_h q$. 由\eqref{eq:chap4-whole-space-localized-stokes}得到
\[
 \left\{
 \begin{aligned}
  -\Delta\delta_h w+\nabla\delta_h q&=\delta_h\widetilde f,\\
  \operatorname{div}\delta_h w&=\delta_h\widetilde g.
 \end{aligned}
 \right.
\]
使用 Nečas 不等式 Theorem~\ref{thm:chap4-necas-inequality}, 得到
\[
 \|\delta_hq\|_{L^2}\leq C\bigl(\|\delta_hq\|_{H^{-1}}+\|\nabla\delta_hq\|_{H^{-1}}\bigr).
\]
由 Lemma~\ref{lem:chap3-translation-sobolev-bound}, 第一项 $\|\delta_hq\|_{H^{-1}}$ 以 $|h|\|q\|_{L^2}$ 为界. 用方程组的第一个方程表示第二项:
\[
 \|\nabla\delta_hq\|_{H^{-1}}
 =\|\delta_h\widetilde f+\Delta\delta_hw\|_{H^{-1}}
 \leq C|h|\|\widetilde f\|_{L^2}+\|\nabla\delta_hw\|_{L^2}.
\]
故最终得到
\begin{equation}
 \|\delta_hq\|_{L^2}
 \leq C|h|\bigl(\|q\|_{L^2}+\|\widetilde f\|_{L^2}\bigr)
 +C\|\nabla\delta_hw\|_{L^2}.
 \label{eq:chap4-localized-pressure-difference}
\end{equation}

现在可将紧支的 $\delta_hw$ 作为第一个方程的测试函数, 并借助第二个方程得到
\[
 \int_{\mathbb R^d}|\nabla\delta_hw|^2\,\mathrm dx
 =\langle\delta_h\widetilde f,\delta_hw\rangle_{H^{-1},H^1}
 +\int_\Omega\delta_hq\,\delta_h\widetilde g\,\mathrm dx
\]
\[
 \leq\|\delta_h\widetilde f\|_{H^{-1}}\|\delta_hw\|_{H^1}
 +\|\delta_hq\|_{L^2}\|\delta_h\widetilde g\|_{L^2}
\]
\[
 \leq C|h|\|\widetilde f\|_{L^2}
 \bigl(\|\nabla\delta_hw\|_{L^2}+|h|\|w\|_{H^1}\bigr)
 +C|h|\|\widetilde g\|_{H^1}\|\delta_hq\|_{L^2}.
\]
利用\eqref{eq:chap4-localized-pressure-difference}, 可得
\[
 \|\nabla\delta_hw\|_{L^2}^2
 \leq C|h|\|\widetilde f\|_{L^2}
 \bigl(\|\nabla\delta_hw\|_{L^2}+|h|\|w\|_{H^1}\bigr)
\]
\[
 \qquad+C|h|^2\|\widetilde g\|_{H^1}
 \bigl(\|q\|_{L^2}+\|\widetilde f\|_{L^2}\bigr)
 +C|h|\|\widetilde g\|_{H^1}\|\nabla\delta_hw\|_{L^2}.
\]
由 Young 不等式, 这给出
\[
 \|\nabla\delta_hw\|_{L^2}^2
 \leq C|h|^2\bigl(\|\widetilde f\|_{L^2}^2+\|w\|_{H^1}^2
 +\|\widetilde g\|_{H^1}^2+\|q\|_{L^2}^2\bigr).
\]
现在使用估计\eqref{eq:chap4-localized-stokes-f-bound}--\eqref{eq:chap4-localized-stokes-q-bound}, 最终得到对任意 $h\in\mathbb R^d$,
\[
 \|\nabla\delta_hw\|_{L^2}^2
 \leq C(\varphi)|h|^2\bigl(\|f\|_{L^2(\omega_2)}^2+\|g\|_{H^1(\omega_2)}^2
 +\|v\|_{H^1(\omega_2)}^2+\|p\|_{L^2(\omega_2)}^2\bigr).
\]
由 Proposition~\ref{prop:chap3-sobolev-difference-quotient-characterization}, 这使我们可以断定 $w\in(H^2(\mathbb R^d))^d$, 并得到所需的 $\|w\|_{H^2(\mathbb R^d)}$ 估计. 此外, 可立即从\eqref{eq:chap4-whole-space-localized-stokes}的第一个方程推出 $\nabla q\in(L^2(\Omega))^d$, 且
\[
 \|w\|_{H^2(\omega_1)}+\|p\|_{H^1(\omega_1)}
 \leq\|w\|_{H^2}+\|q\|_{H^1}
\]
\[
 \leq C(\varphi)\bigl(\|f\|_{L^2(\omega_2)}+\|g\|_{H^1(\omega_2)}
 +\|v\|_{H^1(\omega_2)}+\|p\|_{L^2(\omega_2)}\bigr).
\]
\end{Proof}

\par\noindent\refstepcounter{chapterfoursectionsixremark}\textbf{Remark~\thechapterfoursectionsixremark:}\label{rem:chap4-stokes-local-exhaustion}\quad
设 $v\in(H^1(\Omega))^d$, $p\in L^2(\Omega)$, $f\in(L^2(\Omega))^d$, 且 $g\in H^1(\Omega)$. 考虑由下式定义的 $\mathbb R^d$ 中的开集族:
\[
 \omega_k=\left\{x\in\Omega,\ d(x,\partial\Omega)>\frac1k\right\}.
\]
对所有充分大的 $k$, 有
\[
 \omega_k\subset\overline{\omega_k}\subset\omega_{k+1}\subset\cdots.
\]
前一定理给出的估计可以具体写成
\[
 \|v\|_{H^2(\omega_k)}^2+\|p\|_{H^1(\omega_k)}^2
 \leq C_k\bigl(\|f\|_{L^2(\Omega)}^2+\|g\|_{H^1(\Omega)}^2
 +\|v\|_{H^1(\Omega)}^2+\|p\|_{L^2(\Omega)}^2\bigr),
\]
其中 $C_k$ 与可取成在 $\omega_k$ 上等于 $1$, 在 $\Omega\setminus\omega_{k+1}$ 上等于 $0$ 的截断函数 $\varphi_k$ 有关. 上面的证明表明
\[
 C_k\sim\|\nabla\varphi_k\|_{L^\infty}^2+\|\Delta\varphi_k\|_{L^\infty}^2.
\]
简单计算表明, 当 $k$ 趋于无穷时, $C_k$ 的行为本质上与 $k^8$ 相同. 这说明上述方法不能在 $k$ 趋于无穷时取极限, 以得到整个 $\Omega$ 上解的正则性.

这尤其是因为椭圆问题的解直到边界的正则性, 如 Stokes 问题中的情形, 强烈依赖于解所满足的边界条件的类型和正则性, 也依赖于区域 $\Omega$ 的正则性.


\paragraph{直到边界的正则性}
\label{subsubsec:chap4-stokes-boundary-regularity}

上一小节证明了 Stokes 问题的局部正则性结果. 为完成正则性结果的证明, 现在需要得到解直到边界的正则性. 如同 Chapter~\ref{chap:sobolev-spaces} 的 Section~\ref{sec:chap3-laplace-problem} 中对 Laplace 问题所做的那样, 分两个主要步骤进行: 首先利用沿切向向量场的平移得到切向正则性, 然后直接由方程推出完全正则性.

先在半空间 $\mathbb R_+^d$ 情形说明这一策略, 因为此时计算简单得多; 然后利用 Chapter~\ref{chap:sobolev-spaces} 的 Section~\ref{sec:chap3-boundary-calculus} 中引入的工具考虑一般情形.

\subparagraph{半空间 $\mathbb R_+^d$}
\label{subsubsec:chap4-stokes-half-space}
\leavevmode\par

本段在半空间 $\mathbb R_+^d$ 的简单几何情形说明一般证明策略. 下一段考虑一般情形.

考虑下列问题:
\begin{equation}
 \left\{
 \begin{aligned}
  -\Delta v+v+\nabla p&=f,&&\text{于 }\mathbb R_+^d,\\
  \operatorname{div}v&=0,&&\text{于 }\mathbb R_+^d,\\
  v&=0,&&\text{于 }\partial\mathbb R_+^d,
 \end{aligned}
 \right.
 \label{eq:chap4-half-space-stokes}
\end{equation}
其中 $f\in(L^2(\mathbb R_+^d))^d$. 由于区域无界, 为了保证算子的强制性, 在动量方程中加入了零阶项 $v$, 从而保证解的存在唯一性. 因此, 该方程组是一个模型问题, 我们用它说明下文在充分光滑有界区域 $\Omega$ 上使用的方法.

由算子的强制性, 已知存在唯一的
\[
 (v,p)\in(H_0^1(\mathbb R_+^d))^d\times L_{\mathrm{loc}}^2(\mathbb R_+^d),
\]
其中压力只确定到一个常数, 它给出前述问题的解并满足
\begin{equation}
 \|v\|_{H_0^1}+\|\nabla p\|_{H^{-1}}\leq C\|f\|_{H^{-1}}.
 \label{eq:chap4-half-space-stokes-energy}
\end{equation}

\par\noindent\refstepcounter{chapterfoursectionsixremark}\textbf{Remark~\thechapterfoursectionsixremark:}\label{rem:chap4-half-space-pressure}\quad
需要注意, 不能断言 $p\in L^2(\mathbb R_+^d)$, 甚至不能断言 $p\in H^{-1}(\mathbb R_+^d)$, 见 Proposition~\ref{prop:chap4-hminusone-poincare}. 特别地, 与有界区域的情形不同, 这里不能考虑平均值为零的压力.

我们要使用 Chapter~\ref{chap:sobolev-spaces} 的 Section~\ref{subsec:chap3-tangential-sobolev-spaces} 中引入的通过平移算子对切向 Sobolev 空间所作的刻画.

在这里考虑的 $\Omega=\mathbb R_+^d$ 情形, 容易看出只需考虑常切向向量场 $\theta(x)=e_i$, $i=1,\ldots,d-1$. 相应流 $h\mapsto\tau^i(h,x)$ 由微分方程
\[
 \left\{
 \begin{aligned}
  \frac{\mathrm d}{\mathrm dh}\tau^i(h,x)&=e_i,\\
  \tau^i(0,x)&=x
 \end{aligned}
 \right.
\]
定义. 当然, 这个方程可立即解得 $\tau^i(h,x)=x+he_i$.

向量 $e_i$, $i=1,\ldots,d-1$, 与 $\mathbb R_+^d$ 的边界相切, 所以这些流满足: 对所有 $x\in\mathbb R_+^d$ 和所有 $h\in\mathbb R$, $\tau^i(h,x)$ 仍属于 $\mathbb R_+^d$.

由于几何非常简单, Proposition~\ref{prop:chap3-curved-translation-properties} 的交换估计在这里没有用. 事实上, 容易验证此时平移算子与微分算子严格交换:
\[
 \frac{\partial(v\circ\tau^i(h,x))}{\partial x_j}
 =\frac{\partial v}{\partial x_j}\circ\tau^i(h,x),
 \quad\forall j\in\{1,\ldots,d\},\quad\forall i\in\{1,\ldots,d-1\}.
\]
如同在 Chapter~\ref{chap:sobolev-spaces} 中所做的那样, 引入与这些平移算子关联的差分算子
\[
 (\delta_h^iv)(x)=v(\tau^i(h,x))-v(x).
\]
于是可采用 Chapter~\ref{chap:sobolev-spaces} 的 Section~\ref{sec:chap3-laplace-problem} 中处理 Laplace 问题时所用的同样方法. 当然, 压力项和无散度约束的存在带来了额外困难, 下文将予以说明. 特别需要注意, 速度场的切向分量和法向分量并不是以相同方式研究的.

\begin{Proposition}
\label{prop:chap4-half-space-stokes-regularity}
假设 $f\in(L^2(\mathbb R_+^d))^d$, 则方程组\eqref{eq:chap4-half-space-stokes}的解 $(v,p)$ 满足 $v\in(H^2(\mathbb R_+^d))^d$ 且 $\nabla p\in L^2(\mathbb R_+^d)$. 此外, 存在 $C>0$ 使得
\[
 \|v\|_{H^2}+\|\nabla p\|_{L^2}\leq C\|f\|_{L^2}.
\]
\end{Proposition}

\begin{Proof}
\textbf{Step 1: 切向正则性.}

设 $\phi\in(H_0^1(\mathbb R_+^d))^d$ 且 $\operatorname{div}\phi=0$. 函数 $\delta_{-h}^i\phi$ 也无散度, 这是由于上述交换性质, 因而可在问题中取它作为测试函数. 使用 Lemma~\ref{lem:chap3-translation-duality}, 得到 $\delta_h^iv$ 所满足的变分问题
\[
 \int_{\mathbb R_+^d}\nabla(\delta_h^iv):\nabla\phi\,\mathrm dx
 +\int_{\mathbb R_+^d}(\delta_h^iv)\cdot\phi\,\mathrm dx
 =\int_{\mathbb R_+^d}(\delta_h^if)\cdot\phi\,\mathrm dx,
\]
它对所有满足 $\operatorname{div}\phi=0$ 的 $\phi\in(H_0^1(\mathbb R_+^d))^d$ 成立. 在该方程中取 $\phi=\delta_h^iv$, 得到
\[
 \|\delta_h^iv\|_{H^1}^2
 =\langle\delta_h^if,\delta_h^iv\rangle_{H^{-1},H_0^1}.
\]
利用假设 $f\in L^2(\mathbb R_+^d)$ 和 Lemma~\ref{lem:chap3-translation-sobolev-bound}, 得到不等式
\[
 \|\delta_h^iv\|_{H^1}\leq|h|\|f\|_{L^2}.
\]
这说明 $v$ 在整个 $\mathbb R_+^d$ 上具有切向正则性. 更准确地说, 像 Proposition~\ref{prop:chap3-sobolev-difference-quotient-characterization} 的证明那样进行, 可推出对每个 $i\in\{1,\ldots,d-1\}$, $\partial v/\partial x_i\in(H^1(\mathbb R_+^d))^d$. 换言之, 除 $i=j=d$ 外, $\partial^2v/\partial x_i\partial x_j$ 对所有 $i,j$ 都属于 $L^2(\mathbb R_+^d)$.

现在可以推出压力的切向正则性. 事实上, 对任意 $i\leq d-1$, 由\eqref{eq:chap4-half-space-stokes}得到
\[
 \frac{\partial}{\partial x_i}(\nabla p)
 =\frac{\partial f}{\partial x_i}-\frac{\partial v}{\partial x_i}
 +\frac{\partial\Delta v}{\partial x_i}
 =\frac{\partial f}{\partial x_i}-\frac{\partial v}{\partial x_i}
 +\operatorname{div}\left(\nabla\frac{\partial v}{\partial x_i}\right).
\]
注意, $\nabla\partial v/\partial x_i$ 项中没有 $\partial^2v/\partial x_d^2$ 型二阶导数, 因此前面的论证表明该项属于 $(L^2(\mathbb R_+^d))^d$. 取该项的散度并使用 $f\in(L^2(\mathbb R_+^d))^d$, 得到
\[
 \nabla\frac{\partial p}{\partial x_i}
 =\frac{\partial}{\partial x_i}(\nabla p)
 \in(H^{-1}(\mathbb R_+^d))^d.
\]
此外, 由\eqref{eq:chap4-half-space-stokes-energy}已知 $\partial p/\partial x_i\in H^{-1}(\mathbb R_+^d)$. 因此由 Nečas 不等式 Theorem~\ref{thm:chap4-necas-inequality}, 已证明对所有 $i\leq d-1$, 分布 $\partial p/\partial x_i$ 属于 $L^2(\mathbb R_+^d)$ 并满足预期估计.

\textbf{Step 2: 完全正则性.}

为得到解的完全正则性, 还需研究 $v$ 和 $p$ 沿 $e_d$ 方向的导数. 与 Laplace 问题一样, 这些性质可直接由偏微分方程以及刚得到的切向正则性推出.

\emph{Substep 1: 速度场法向分量的完全正则性.}

将散度方程对 $x_d$ 求导, 得到
\[
 \frac{\partial^2v_d}{\partial x_d^2}
 =-\sum_{i=1}^{d-1}\frac{\partial^2v_i}{\partial x_i\partial x_d}
 \in L^2(\mathbb R_+^d).
\]
因此已经证明法向分量 $v_d\in H^2(\mathbb R_+^d)$.

\emph{Substep 2: 压力的完全正则性.}

将 Stokes 问题中动量方程的第 $d$ 个分量写为
\[
 \frac{\partial p}{\partial x_d}=\Delta v_d-v_d+f_d.
\]
利用上面证明的 $v_d$ 的正则性, 立即看出 $\partial p/\partial x_d\in L^2(\mathbb R_+^d)$. 因此已经证明 $p\in H^1(\mathbb R_+^d)$.

\emph{Substep 3: 速度场切向分量的完全正则性.}

对 $i\leq d-1$, 可将方程组的第 $i$ 个方程写为
\[
 \frac{\partial^2v_i}{\partial x_d^2}
 =-\sum_{j=1}^{d-1}\frac{\partial^2v_i}{\partial x_j^2}
 +v_i+\frac{\partial p}{\partial x_i}-f_i.
\]
由此推出 $\partial^2v_i/\partial x_d^2\in L^2(\mathbb R_+^d)$, 从而 $v_i\in H^2(\mathbb R_+^d)$.

于是正则性定理得到证明. 更准确地说, 已建立某个 $C>0$ 使得
\[
 \|v\|_{H^2}+\|\nabla p\|_{L^2}
 \leq C\bigl(\|f\|_{L^2}+\|\nabla p\|_{H^{-1}}\bigr).
\]
下文将在一般光滑有界区域的情形遵循同一方案.
\end{Proof}

\subparagraph{切向正则性}
\label{subsubsec:chap4-stokes-tangential-regularity}
\leavevmode\par

我们回忆, 已在 Theorem~\ref{thm:chap4-stokes-local-regularity} 中证明 Stokes 问题的解具有局部正则性 $(H_{\mathrm{loc}}^2(\Omega))^d\times H_{\mathrm{loc}}^1(\Omega)$. 现在证明, 若数据直到边界都是正则的, 则同一解具有切向正则性 $(H_{\mathrm{tang}}^2(\Omega))^d\times H_{\mathrm{tang}}^1(\Omega)$.

这里使用 Chapter~\ref{chap:sobolev-spaces} 的 Section~\ref{sec:chap3-boundary-calculus} 中引入的记号.

\begin{Theorem}
\label{thm:chap4-stokes-tangential-regularity}
设 $\Omega$ 是 $\mathbb R^d$ 中 $C^{1,1}$ 类的有界连通区域, $f\in(L^2(\Omega))^d$, 且 $g\in H^1(\Omega)\cap L_0^2(\Omega)$. 若 $(v,p)$ 是下列问题在 $(H_0^1(\Omega))^d\times L_0^2(\Omega)$ 中的唯一解:
\[
 \left\{
 \begin{aligned}
  -\Delta v+\nabla p&=f,&&\text{于 }\Omega,\\
  \operatorname{div}v&=g,&&\text{于 }\Omega,\\
  v&=0,&&\text{于 }\partial\Omega,
 \end{aligned}
 \right.
\]
则 $v\in(H_{\mathrm{tang}}^2(\Omega))^d$, $p\in H_{\mathrm{tang}}^1(\Omega)$, 且对任意满足\eqref{eq:chap3-admissible-tangential-field}中 $k=0$ 情形的向量场 $\theta$, 存在 $C(\Omega,\theta)>0$ 使得
\[
 |||v|||_{\theta,2}+|||p|||_{\theta,1}
 \leq C(\Omega,\theta)\bigl(\|f\|_{L^2}+\|g\|_{H^1}\bigr).
\]
\end{Theorem}

这个结果的证明需要若干步骤, 它们受到半空间情形中计算的启发, 见 Section~\ref{subsubsec:chap4-stokes-half-space}. 首先求差分 $\delta_h^\theta v=v\circ\tau^\theta(h,\cdot)-v$ 和 $\delta_h^\theta p=p\circ\tau^\theta(h,\cdot)-p$ 所满足的方程, 并注意微分运算.

下面的 Lemma 给出微分算子与沿切向场的平移算子之间交换子的适当表达式. 在前面考虑的半空间情形, 由于切向平移的结构特别简单, 这些交换子为零.

\begin{Lemma}
\label{lem:chap4-stokes-tangential-commutators}
有如下等式:
\begin{equation}
 \left\{
 \begin{aligned}
  -\Delta(\delta_h^\theta v)
  &=-\delta_h^\theta(\Delta v)-[\operatorname{div},\tau_h^\theta]\nabla v
    -\operatorname{div}([\nabla,\tau_h^\theta]v),\\
  \nabla\delta_h^\theta p
  &=\delta_h^\theta\nabla p+[\nabla,\tau_h^\theta]p,\\
  \operatorname{div}(\delta_h^\theta v)
  &=\delta_h^\theta(\operatorname{div}v)+[\operatorname{div},\tau_h^\theta]v.
 \end{aligned}
 \right.
 \label{eq:chap4-stokes-tangential-commutators}
\end{equation}
\end{Lemma}

\begin{Proof}
(Lemma~\ref{lem:chap4-stokes-tangential-commutators}).
这只是简单的纯代数计算.
\end{Proof}

\begin{Proof}
(Theorem~\ref{thm:chap4-stokes-tangential-regularity}).
设 $h\in(0,1)$, 并令 $\theta$ 为与 $\Omega$ 的边界相切的向量场, 即满足\eqref{eq:chap3-admissible-tangential-field}中 $k=0$ 的情形. 记
\[
 w=\delta_h^\theta v,\qquad\pi=\delta_h^\theta p.
\]
利用\eqref{eq:chap4-stokes-tangential-commutators}, 可见 $(w,\pi)$ 满足方程组
\begin{equation}
 \left\{
 \begin{aligned}
  -\Delta w+\nabla\pi&=F,&&\text{于 }\Omega,\\
  \operatorname{div}w&=G,&&\text{于 }\Omega,
 \end{aligned}
 \right.
 \label{eq:chap4-tangential-difference-stokes}
\end{equation}
其中
\begin{equation}
 F=\delta_h^\theta f+[\nabla,\tau_h^\theta]p
 -[\operatorname{div},\tau_h^\theta]\nabla v
 -\operatorname{div}([\nabla,\tau_h^\theta]v),
 \label{eq:chap4-tangential-difference-force}
\end{equation}
以及
\begin{equation}
 G=\delta_h^\theta g+\delta_h^\theta(\operatorname{div}v)
 +[\operatorname{div},\tau_h^\theta]v.
 \label{eq:chap4-tangential-difference-divergence}
\end{equation}
注意, 对固定的 $h$, $\tau^\theta(h,\cdot)$ 是 $\partial\Omega$ 到自身的双射, 因而 $w\in(H_0^1(\Omega))^d$. 现在估计 $F$ 在 $(H^{-1}(\Omega))^d$ 中的范数. 由\eqref{eq:chap3-curved-translation-gradient-commutator}和\eqref{eq:chap3-curved-translation-seminorm-bound}, 得到
\[
 \|\delta_h^\theta f\|_{H^{-1}}\leq Ch\|f\|_{L^2},
\]
\[
 \|[\operatorname{div},\tau_h^\theta]\nabla v\|_{H^{-1}}
 \leq Ch\|\nabla v\|_{L^2},
 \qquad
 \|\operatorname{div}([\nabla,\tau_h^\theta]v)\|_{H^{-1}}
 \leq\|[\nabla,\tau_h^\theta]v\|_{L^2}
 \leq Ch\|v\|_{H^1}.
\]
再用 Theorem~\ref{thm:chap4-stokes-nonhomogeneous-wellposedness} 给出的初始问题解在 $(H_0^1(\Omega))^d$ 中的估计, 得到
\[
 \|F\|_{H^{-1}}\leq Ch\bigl(\|f\|_{L^2}+\|g\|_{H^1}\bigr).
\]
同理, 可在 $L^2(\Omega)$ 中估计 $G$:
\[
 \|G\|_{L^2}\leq Ch\|g\|_{H^1}+Ch\|v\|_{H^1}
 \leq Ch\bigl(\|f\|_{L^2}+\|g\|_{H^1}\bigr).
\]
接着对问题\eqref{eq:chap4-tangential-difference-stokes}作经典能量估计, 得到
\begin{align}
 \|\nabla w\|_{L^2}^2
 &\leq\|F\|_{H^{-1}}\|\nabla w\|_{L^2}
 +\left|\int_\Omega\pi G\,\mathrm dx\right|\notag\\
 &\leq Ch\bigl(\|f\|_{L^2}+\|g\|_{H^1}\bigr)\|\nabla w\|_{L^2}
 +\|\pi\|_{L^2}\|G\|_{L^2}\notag\\
 &\leq Ch\bigl(\|f\|_{L^2}+\|g\|_{H^1}\bigr)\|\nabla w\|_{L^2}
 +Ch\|\pi\|_{L^2}\bigl(\|f\|_{L^2}+\|g\|_{H^1}\bigr).
 \label{eq:chap4-tangential-energy-estimate}
\end{align}
此外, 从原文所称的\MBUnresolvedReference{eq:source-iv-6-4}{Equation (IV.6.4)}直接读出
\[
 \|\nabla\pi\|_{H^{-1}}
 \leq\|\Delta w\|_{H^{-1}}+\|F\|_{H^{-1}}
 \leq C\|\nabla w\|_{L^2}+Ch\bigl(\|f\|_{L^2}+\|g\|_{H^1}\bigr),
\]
并由\eqref{eq:chap3-curved-translation-seminorm-bound}和 Theorem~\ref{thm:chap4-stokes-nonhomogeneous-wellposedness} 有
\[
 \|\pi\|_{H^{-1}}=\|\delta_h^\theta p\|_{H^{-1}}
 \leq Ch\|p\|_{L^2}
 \leq Ch\bigl(\|f\|_{L^2}+\|g\|_{H^1}\bigr).
\]
最后由 Nečas 不等式 Theorem~\ref{thm:chap4-necas-inequality}, 得到
\begin{equation}
 \|\pi\|_{L^2}
 \leq C\bigl(\|\pi\|_{H^{-1}}+\|\nabla\pi\|_{H^{-1}}\bigr)
 \leq C\|\nabla w\|_{L^2}+Ch\bigl(\|f\|_{L^2}+\|g\|_{H^1}\bigr).
 \label{eq:chap4-tangential-pressure-estimate}
\end{equation}
因此, 由不等式\eqref{eq:chap4-tangential-energy-estimate}和\eqref{eq:chap4-tangential-pressure-estimate}及 Young 不等式,
\[
 \|\nabla w\|_{L^2}^2
 \leq Ch\bigl(\|f\|_{L^2}+\|g\|_{H^1}\bigr)\|\nabla w\|_{L^2}
 +Ch^2\bigl(\|f\|_{L^2}^2+\|g\|_{H^1}^2\bigr)
\]
\[
 \leq Ch^2\bigl(\|f\|_{L^2}^2+\|g\|_{H^1}^2\bigr)
 +\frac12\|\nabla w\|_{L^2}^2.
\]
这使我们可以结合\eqref{eq:chap4-tangential-pressure-estimate}, 得到不等式
\[
 \|\nabla\delta_h^\theta v\|_{L^2}^2+\|\delta_h^\theta p\|_{L^2}^2
 \leq Ch^2\bigl(\|f\|_{L^2}^2+\|g\|_{H^1}^2\bigr),
\]
这恰好表明
\[
 |||v|||_{\theta,2}^2+|||p|||_{\theta,1}^2
 \leq C\bigl(\|f\|_{L^2}^2+\|g\|_{H^1}^2\bigr).
\]
\end{Proof}

\subparagraph{直到边界的完全正则性}
\label{subsubsec:chap4-stokes-complete-boundary-regularity}
\leavevmode\par

现在使用 Chapter~\ref{chap:sobolev-spaces} 中建立的 $\Omega$ 边界管状邻域 $O_{s_0}^\rho$ 内的坐标系. 刚才已看到, 在 Theorem~\ref{thm:chap4-stokes-regularity} 的 $k=0$ 假设下, Stokes 问题的解属于 $(H_{\mathrm{tang}}^2(\Omega))^d\times H_{\mathrm{tang}}^1(\Omega)$. 为证明解在整个 $\Omega$ 上的正则性, 只需证明它在 $O_{s_0}^\rho$ 中的正则性. 切向/法向坐标系使切向变量和法向变量的作用能够分离, 正如 Section~\ref{subsubsec:chap4-stokes-half-space} 中对简单得多的半空间 $\mathbb R_+^d$ 所做的那样.

回忆, 在 Chapter~\ref{chap:sobolev-spaces} 的 Section~\ref{sec:chap3-boundary-calculus} 中, 我们用 $\langle x,y\rangle$ 表示两个向量 $x,y$ 的 Euclidean 标量积, 而下文将该标量积记为 $x\cdot y$.

在开集 $O_{s_0}^\rho$ 中, $v$ 的散度方程写为
\begin{equation}
 \operatorname{Tr}(\nabla_Tv)+\nabla_Nv_N+v\cdot\widetilde G
 =g\in H^1(O_{s_0}^\rho).
 \label{eq:chap4-stokes-boundary-divergence}
\end{equation}
利用前面证明的 $v$ 的切向正则性, 可见该等式第一项属于 $H^1(O_{s_0}^\rho)$. 此外,
\begin{equation}
 \nabla_N(v\cdot\widetilde G)
 =(\nabla_Nv)\cdot\widetilde G+v\cdot\nabla_N\widetilde G.
 \label{eq:chap4-stokes-normal-product}
\end{equation}
由于 $\widetilde G$ 由 $|\ln\rho|$ 控制, 且 $\nabla v\in(L^2(\Omega))^{d\times d}$, 推出
\begin{equation}
 (\nabla_Nv)\cdot\widetilde G\in L^q(\Omega),\qquad\forall q<2.
 \label{eq:chap4-stokes-normal-product-first}
\end{equation}
利用 $|\nabla_N\widetilde G|\leq C/\rho$ 以及 $v=0$ 于 $\partial\Omega$, Hardy 不等式说明\eqref{eq:chap4-stokes-normal-product}右端第二项属于 $L^2(\Omega)$, 并满足
\begin{equation}
 \|v\cdot(\nabla_N\widetilde G)\|_{L^2}\leq C\|v\|_{H^1}.
 \label{eq:chap4-stokes-normal-product-second}
\end{equation}
因此由\eqref{eq:chap4-stokes-boundary-divergence}和\eqref{eq:chap4-stokes-normal-product}, 有
\[
 \nabla_N\nabla_Nv_N\in L^q(O_{s_0}^\rho),\qquad\forall q<2.
\]

令 $q<2$, 并证明 $v_N\in W^{2,q}(O_{s_0}^\rho)$. 首先证明 $\nabla_Nv_N\in(W^{1,q}(\Omega))^d$. 为此计算
\[
 \nabla(\nabla_Nv_N)=\nabla_T(\nabla_Nv_N)+(\nabla_N\nabla_Nv_N)\nu
\]
\[
 =\nabla_T\bigl((\nabla_Nv)\cdot\nu\bigr)
 +\nabla_T\bigl(v\cdot(\nabla_N\nu)\bigr)
 +(\nabla_N\nabla_Nv_N)\nu,
\]
其中第一项也可写为 $\nabla_T\bigl((\nabla v)\cdot\nu\cdot\nu\bigr)$.
刚才已证明最后一项属于 $L^q$. 由于 $\nu$ 是 Lipschitz 连续的且 $\nabla v\in(H_{\mathrm{tang}}^1(\Omega))^{d\times d}$, 第一项属于 $L^2\subset L^q$. 最后, 利用 $\nu$ 的性质和 Hardy 不等式, 可见第二项满足
\begin{equation}
 \nabla_T\bigl(v\cdot(\nabla_N\nu)\bigr)\in L^2(\Omega)\subset L^q(\Omega).
 \label{eq:chap4-stokes-normal-velocity-second-term}
\end{equation}
这实际上证明了 $\nabla_Nv_N\in(W^{1,q}(\Omega))^d$.

最后写成
\[
 \nabla v_N=\nabla_Tv_N+(\nabla_Nv_N)\nu.
\]
由于每一项都已证明属于 $(W^{1,q}(\Omega))^d$, 推出 $v_N\in W^{2,q}(\Omega)$.

将 $\Delta v$ 的法向部分计算为
\[
 (\Delta v)_N=(\Delta v)\cdot\nu
 =\Delta(v_N)-v\cdot(\Delta\nu)-2(\nabla v):(\nabla\nu).
\]
作与上面相同的计算. 特别地, 利用 $|\Delta\nu|\sim1/\rho$ 和 Hardy 不等式, 证明第二项满足
\begin{equation}
 v\cdot(\Delta\nu)\in L^2(\Omega),
 \label{eq:chap4-stokes-laplacian-normal-term}
\end{equation}
再利用 $v_N\in W^{2,q}(\Omega)$. 由此推出 $(\Delta v)_N\in L^q(\Omega)$. 动量方程投影到法向方向给出
\[
 -(\Delta v)_N+\nabla_Np=f_N,
\]
从而
\[
 \nabla_Np\in L^q(\Omega).
\]
由于已经知道 $\nabla_Tp\in(L^2(\Omega))^d$, 最终得到 $p\in W^{1,q}(\Omega)$.

令 $1\leq i\leq d$, 并考察动量方程的第 $i$ 个分量:
\[
 -\operatorname{Tr}(\nabla_T\nabla v_i)-\nabla_N\nabla_Nv_i
 -(\nabla v_i)\cdot\widetilde G+\frac{\partial p}{\partial x_i}=f_i.
\]
除第二项外, 所有项都已知属于 $L^q(O_{s_0}^\rho)$, 因此 $\nabla_N\nabla_Nv_i\in L^q(O_{s_0}^\rho)$. 像前面一样可证明 $\nabla_Nv_i\in W^{1,q}(O_{s_0}^\rho)$. 事实上,
\[
 \nabla(\nabla_Nv_i)=\nabla_T(\nabla_Nv_i)+(\nabla_N\nabla_Nv_i)\nu
 =\nabla_T\bigl((\nabla v)\cdot\nu\cdot e_i\bigr)+(\nabla_N\nabla_Nv_i)\nu,
\]
利用 $\nu$ 的性质以及上面得到的 $v$ 的全部正则性, 每一项都属于 $L^q$.

为结束证明, 写成
\[
 \nabla v_i=\nabla_Tv_i+(\nabla_Nv_i)\nu.
\]
于是 $\nabla v_i\in(W^{1,q}(O_{s_0}^\rho))^d$, 从而 $v\in(W^{2,q}(\Omega))^d$.

我们已经证明对所有 $q<2$, $v\in(W^{2,q}(\Omega))^d$. 由 Sobolev 嵌入, 推出存在某个 $r>2$ 使得 $\nabla v\in(L^r(\Omega))^{d\times d}$. 结果, 
\eqref{eq:chap4-stokes-normal-product-first}对 $q=2$ 也成立, 上述证明的后续步骤在此情形仍然成立. 最终得到 $v\in(H^2(\Omega))^d$.

至此完成了 Theorem~\ref{thm:chap4-stokes-regularity} 在 $k=0$ 情形的证明.

\subsubsection{高阶正则性}
\label{subsec:chap4-stokes-higher-regularity}

现在转向 Stokes 问题椭圆正则性结果的一般情形.

\begin{Proof}
(Theorem~\ref{thm:chap4-stokes-regularity}).
设 $k\geq0$, $\Omega$ 是 $\mathbb R^d$ 中 $C^{k+1,1}$ 类的有界连通区域, $f\in(H^k(\Omega))^d$, 且 $g\in H^{k+1}(\Omega)$. 我们要证明下列问题的解 $(v,p)\in(H_0^1(\Omega))^d\times L_0^2(\Omega)$:
\[
 \left\{
 \begin{aligned}
  -\Delta v+\nabla p&=f,&&\text{于 }\Omega,\\
  \operatorname{div}v&=g,&&\text{于 }\Omega,\\
  v&=0,&&\text{于 }\partial\Omega,
 \end{aligned}
 \right.
\]
属于 $(H^{k+2}(\Omega))^d\times H^{k+1}(\Omega)$. 为此使用归纳法.

$k=0$ 的情形已在 Section~\ref{subsec:chap4-stokes-first-regularity} 中证明. 取 $k\geq1$, 并假设把 $k$ 换为 $k-1$ 时已经证明了结果:
\[
 \|v\|_{H^{k+1}}+\|p\|_{H^k}
 \leq C\bigl(\|f\|_{H^{k-1}}+\|g\|_{H^k}\bigr).
\]
可以像以前一样, 首先把问题化到全空间 $\mathbb R^d$, 然后沿坐标轴作平移, 以证明局部正则性. 为得到切向正则性, 再次采用 Theorem~\ref{thm:chap4-stokes-tangential-regularity} 的证明中的记号. 回忆, 已经证明 $w=\delta_h^\theta v$ 和 $\pi=\delta_h^\theta p$ 是 Stokes 问题\eqref{eq:chap4-tangential-difference-stokes}的解, 其第二端由\eqref{eq:chap4-tangential-difference-force}和\eqref{eq:chap4-tangential-difference-divergence}定义.

利用应用于初始 Stokes 问题的归纳假设, $f$ 和 $g$ 的正则性以及\eqref{eq:chap3-curved-translation-gradient-commutator}和\eqref{eq:chap3-curved-translation-seminorm-bound}的性质, 立即得到
\[
 \|F\|_{H^{k-1}}\leq Ch\bigl(\|f\|_{H^k}+\|g\|_{H^{k+1}}\bigr),
\]
\[
 \|G\|_{H^k}\leq Ch\bigl(\|f\|_{H^k}+\|g\|_{H^{k+1}}\bigr).
\]
因此, 对 Stokes 问题\eqref{eq:chap4-tangential-difference-stokes}应用归纳假设, 得到
\[
 \|\delta_h^\theta v\|_{H^{k+1}}^2+\|\delta_h^\theta p\|_{H^k}^2
 \leq Ch^2\bigl(\|f\|_{H^k}^2+\|g\|_{H^{k+1}}^2\bigr).
\]
由切向 Sobolev 空间的刻画, 即原文所称的\MBUnresolvedReference{thm:source-iii-2-32-higher-stokes}{Theorem III.2.32}, 这就给出结果.

最后, 再次使用 $\Omega$ 内 $\partial\Omega$ 的管状邻域 $O_{s_0}^\rho$ 中的切向和法向坐标, 从切向正则性推出解直到边界的正则性. 同样, 这个证明要求区域 $\Omega$ 具有高一阶正则性.
\end{Proof}

\subsubsection{Stokes 问题的 $L^q$ 理论}
\label{subsec:chap4-stokes-lq-theory}

本节的最后一个小节旨在非常简略地介绍在 $L^q(\Omega)$ 空间中提出的 Stokes 问题理论的若干内容. 关于这一问题的完整研究, 参见\MBUnresolvedReference{bib:stokes-lq-monograph}{[63]}.

首先陈述由 Agmon, Douglis 和 Nirenberg 证明的下列定理. 它假设 $(W^{2,q}(\Omega))^d\times W^{1,q}(\Omega)$ 中存在一个解, 并由此得到该解的更高正则性. 事实上, 这个定理是\MBUnresolvedReference{bib:agmon-douglis-nirenberg-2}{[2]}和\MBUnresolvedReference{bib:agmon-douglis-nirenberg-3}{[3]}中所建立的更一般定理的一个特例.

我们强调, 这个定理并不是光滑解的存在性定理.

\begin{Theorem}{Agmon--Douglis--Nirenberg}
\label{thm:chap4-stokes-agmon-douglis-nirenberg}
设 $\Omega$ 是 $\mathbb R^d$ 中 $C^{k+2}$ 类的有界区域, 其中 $k$ 是正整数. 设 $1<q<+\infty$, 并设数据
\[
 f\in(L^q(\Omega))^d,\qquad g\in W^{1,q}(\Omega),\qquad
 v_b\in(W^{2-1/q,q}(\partial\Omega))^d
\]
满足相容性条件
\[
 \int_\Omega g\,\mathrm dx=\int_{\partial\Omega}v_b\cdot\nu\,\mathrm d\sigma.
\]
假设存在一对
\[
 (v,p)\in(W^{2,q}(\Omega))^d\times\bigl(W^{1,q}(\Omega)\cap L_0^2(\Omega)\bigr),
\]
它是 Stokes 问题
\[
 \left\{
 \begin{aligned}
  -\Delta v+\nabla p&=f,&&\text{于 }\Omega,\\
  \operatorname{div}v&=g,&&\text{于 }\Omega,\\
  v&=v_b,&&\text{于 }\partial\Omega
 \end{aligned}
 \right.
\]
的解. 若数据 $(f,g,v_b)$ 具有更高正则性
\[
 (f,g,v_b)\in(W^{k,q}(\Omega))^d\times W^{k+1,q}(\Omega)
 \times(W^{k+2-1/q,q}(\partial\Omega))^d,
\]
则上述问题的解满足
\begin{equation}
 (v,p)\in(W^{k+2,q}(\Omega))^d\times W^{k+1,q}(\Omega),
 \label{eq:chap4-stokes-lq-regularity}
\end{equation}
并有估计
\begin{equation}
 \|v\|_{W^{k+2,q}}+\|p\|_{W^{k+1,q}}
 \leq C\bigl(\|f\|_{W^{k,q}}+\|g\|_{W^{k+1,q}}
 +\|v_b\|_{W^{k+2-1/q,q}}+d_q\|v\|_{L^q}\bigr),
 \label{eq:chap4-stokes-lq-estimate}
\end{equation}
其中 $C$ 只依赖于 $(q,k,\Omega)$, 并且当 $q\geq2$ 时 $d_q=0$, 否则 $d_q=1$.
\end{Theorem}

这个定理是真正的正则性定理, 因为它假设已经知道解的存在性. 为得到兼具存在性与正则性的结果, 例如需要引用 Cattabriga 在\MBUnresolvedReference{bib:cattabriga-stokes}{[39]}中的结果, 另见\MBUnresolvedReference{bib:stokes-lq-reference}{[7]}. 该结果由作者在三维情形证明, 这里不加证明地给出.

\begin{Theorem}{Cattabriga}
\label{thm:chap4-stokes-cattabriga}
设 $k$ 是满足 $k\geq-1$ 的整数, $\Omega$ 是 $\mathbb R^3$ 中 $C^m$ 类的有界连通区域, 其中 $m=\max(2,k+2)$. 对所有满足
\[
 (f,g,v_b)\in(W^{k,q}(\Omega))^d\times W^{k+1,q}(\Omega)
 \times(W^{k+2-1/q,q}(\partial\Omega))^d,
\]
且以 $1<q<+\infty$ 满足相容性条件
\[
 \int_\Omega g\,\mathrm dx=\int_{\partial\Omega}v_b\cdot\nu\,\mathrm d\sigma
\]
的数据, 存在唯一的下列问题的解 $(v,p)$:
\[
 \left\{
 \begin{aligned}
  -\Delta v+\nabla p&=f,&&\text{于 }\Omega,\\
  \operatorname{div}v&=g,&&\text{于 }\Omega,\\
  v&=v_b,&&\text{于 }\partial\Omega,
 \end{aligned}
 \right.
\]
该解满足\eqref{eq:chap4-stokes-lq-regularity}以及 $d_q=0$ 对所有 $q$ 都成立的估计\eqref{eq:chap4-stokes-lq-estimate}.
\end{Theorem}

注意, 该定理在二维情形也成立, 例如见\MBUnresolvedReference{bib:stokes-lq-two-dimensional}{[122]}.

\subsubsection{div/curl 问题的正则性}
\label{subsec:chap4-divcurl-regularity-proof}

现在可以证明 Section~\ref{subsec:chap4-div-curl-problem} 中陈述的 div/curl 方程组正则性定理.

\begin{Proof}
(Theorem~\ref{thm:chap4-div-curl-normal-boundary}).
假设 $u\in(H^1(\Omega))^3$ 满足 $\operatorname{div}u\in H^1(\Omega)$, $\operatorname{curl}u\in(H^1(\Omega))^3$, 且 $\gamma_\nu(u)\in H^{3/2}(\partial\Omega)$. 我们要证明 $u\in(H^2(\Omega))^3$.

首先注意, 通过使用一个适当的迹提升, 只需考虑 $\gamma_\nu(u)=0$ 的情形.

首先证明解的局部正则性. 设 $\varphi\in\mathcal D(\Omega)$ 并令 $v=\varphi u$. 利用\MBUnresolvedReference{eq:appendix-vector-a4}{(A.4)}和\MBUnresolvedReference{eq:appendix-vector-a5}{(A.5)}, 以及 $v$ 在 $\Omega$ 中紧支的事实, 得到
\[
 \operatorname{div}v=\varphi(\operatorname{div}u)+u\cdot\nabla\varphi\in H^1(\Omega),
\]
\[
 \operatorname{curl}v=\varphi(\operatorname{curl}u)+\nabla\varphi\wedge u\in H^1(\Omega).
\]
因此对任意 $1\leq i\leq3$, $w_i=\partial_{x_i}v$ 满足
\[
 w_i\in(L^2(\Omega))^3,\quad
 \operatorname{div}w_i\in L^2(\Omega),\quad
 \operatorname{curl}w_i\in(L^2(\Omega))^3,\quad
 \gamma_\nu w_i=0.
\]
使用 Theorem~\ref{thm:chap4-hdivcurlnu-equals-h1}, 推出 $w_i\in(H^1(\Omega))^3$, 且
\[
 \|w_i\|_{H^1}\leq C_\varphi
 \bigl(\|u\|_{H^1}+\|\operatorname{div}u\|_{H^1}+\|\operatorname{curl}u\|_{H^1}\bigr).
\]
这对任意 $1\leq i\leq3$ 成立, 因而对任意 $\varphi\in\mathcal D(\Omega)$, $\varphi u=v\in(H^2(\Omega))^3$. 这就证明 $u\in(H_{\mathrm{loc}}^2(\Omega))^3$.

现在只需研究边界 $\partial\Omega$ 的一个邻域中的正则性. 使用 Section~\ref{subsubsec:chap4-stokes-boundary-regularity} 中的同一策略, 先证明 $u$ 的切向正则性. 为此, 考虑一个与 $\Omega$ 边界相切的向量场 $\theta$, 即满足\eqref{eq:chap3-admissible-tangential-field}, 并使用相应流 $(h,x)\mapsto\tau^\theta(h,x)$, 见\eqref{eq:chap3-tangential-flow}. 由于 $\theta$ 足够光滑, 可见 $\delta_h^\theta u=\tau_h^\theta u-u\in(H^1(\Omega))^3$, 并且还满足, 见 Lemma~\ref{lem:chap4-stokes-tangential-commutators},
\[
 \operatorname{div}(\delta_h^\theta u)
 =\delta_h^\theta\operatorname{div}u+[\operatorname{div},\tau_h^\theta]u,
\]
\[
 \operatorname{curl}(\delta_h^\theta u)
 =\delta_h^\theta\operatorname{curl}u+[\operatorname{curl},\tau_h^\theta]u,
\]
\[
 \gamma_\nu(\delta_h^\theta u)
 =\tau_h^\theta u\cdot\nu-u\cdot\nu
 =-(\tau_h^\theta u)\cdot(\tau_h^\theta\nu)+\tau_h^\theta u\cdot\nu
 =-\tau_h^\theta u\cdot(\tau_h^\theta\nu-\nu)
 =-\tau_h^\theta u\cdot\delta_h^\theta\nu.
\]
利用 Proposition~\ref{prop:chap3-curved-translation-properties} 中的交换子估计, 得到
\[
 \|\operatorname{div}(\delta_h^\theta u)\|_{L^2}
 \leq C_\theta|h|\bigl(\|\operatorname{div}u\|_{H^1}+\|u\|_{H^1}\bigr),
\]
\[
 \|\operatorname{curl}(\delta_h^\theta u)\|_{L^2}
 \leq C_\theta|h|\bigl(\|\operatorname{curl}u\|_{H^1}+\|u\|_{H^1}\bigr),
\]
\begin{equation}
 \|\gamma_\nu(\delta_h^\theta u)\|_{H^{1/2}}
 \leq C_\theta|h|\|\nu\|_{W^{2,\infty}}\|u\|_{H^1}.
 \label{eq:chap4-divcurl-tangential-trace-estimate}
\end{equation}
使用 Theorem~\ref{thm:chap4-hdivcurlnu-equals-h1}, 推出 $\delta_h^\theta u$ 的 $H^1$ 范数满足
\[
 \|\delta_h^\theta u\|_{H^1}
 \leq C_{\theta,\nu}|h|
 \bigl(\|u\|_{H^1}+\|\operatorname{div}u\|_{H^1}+\|\operatorname{curl}u\|_{H^1}\bigr).
\]
由 Theorem~\ref{thm:chap3-tangential-space-characterization}, 推出 $u\in(H_{\mathrm{tang}}^2(\Omega))^3$, 并有相应的范数估计.

使用 Proposition~\ref{prop:chap3-divergence-normal-tangential} 得到 $u_N=u\cdot\nu$ 满足
\[
 \nabla_Nu_N=\operatorname{div}u-\operatorname{div}_T u-\frac{\nabla_NJ}{J}u_N,
 \qquad\text{于 }O_{s_0}^\rho.
\]
由于 $\Omega$ 是 $C^{2,1}$ 类的, 已知 $J$ 在 $O_{s_0}^\rho$ 上属于 $C^{1,1}$, 见 Proposition~\ref{prop:chap3-level-surface-jacobian}. 因此, 由假设 $\operatorname{div}u\in H^1(\Omega)$ 及刚才证明的 $u\in(H_{\mathrm{tang}}^2(\Omega))^3$, 推出
\[
 \nabla_N(u\cdot\nu)\in H^1(O_{s_0}^\rho).
\]
再利用已经证明的 $u$ 的切向正则性, 得到 $u\cdot\nu\in H^2(\Omega)$.

最后, 对任意 $1\leq i\leq3$, 令 $w_i=\partial_{x_i}u$, 可见
\[
 w_i\cdot\nu=\partial_{x_i}(u\cdot\nu)-u\cdot(\partial_{x_i}\nu)\in H^1(\Omega),
\]
因为 $\nu\in C^{1,1}(\Omega)$. 因此迹 $\gamma_\nu(w_i)\in H^{1/2}(\partial\Omega)$. 此外,
\[
 \operatorname{div}w_i=\partial_{x_i}(\operatorname{div}u)\in L^2(\Omega),
 \qquad
 \operatorname{curl}w_i=\partial_{x_i}(\operatorname{curl}u)\in(L^2(\Omega))^3.
\]
再次使用 Theorem~\ref{thm:chap4-hdivcurlnu-equals-h1}, 得到 $w_i\in(H^1(\Omega))^3$. 这对所有 $i$ 都成立, 因而得到 $u\in(H^2(\Omega))^3$ 及所需估计.
\end{Proof}

